
\documentclass{article}
\usepackage{colortbl}
\usepackage{makecell}
\usepackage{multirow}
\usepackage{supertabular}

\begin{document}

\newcounter{utterance}

\centering \large Interaction Transcript for game `hot\_air\_balloon', experiment `air\_balloon\_survival\_en\_complexity\_easy', episode 4 with claude{-}sonnet{-}4{-}5{-}20250929{-}t1.0.
\vspace{24pt}

{ \footnotesize  \setcounter{utterance}{1}
\setlength{\tabcolsep}{0pt}
\begin{supertabular}{c@{$\;$}|p{.15\linewidth}@{}p{.15\linewidth}p{.15\linewidth}p{.15\linewidth}p{.15\linewidth}p{.15\linewidth}}
    \# & $\;$A & \multicolumn{4}{c}{Game Master} & $\;\:$B\\
    \hline

    \theutterance \stepcounter{utterance}  
    & & \multicolumn{4}{p{0.6\linewidth}}{
        \cellcolor[rgb]{0.9,0.9,0.9}{
            \makecell[{{p{\linewidth}}}]{
                \texttt{\tiny{[P1$\langle$GM]}}
                \texttt{You are participating in a collaborative negotiation game.} \\
\\ 
\texttt{Together with another participant, you must agree on a single set of items that will be kept. Each of you has your own view of how much each item matters to you (importance). You do not know how the other participant values the items. Additionally, you are given the effort each item demands.} \\
\texttt{You may only agree on a set if the total effort of the selected items does not exceed a shared limit:} \\
\\ 
\texttt{LIMIT = 9612} \\
\\ 
\texttt{Here are the individual item effort values:} \\
\\ 
\texttt{Item effort = \{"B21": 935, "B00": 928, "A08": 725, "A27": 218, "B94": 839, "B53": 970, "A23": 394, "C65": 677, "A89": 946, "B67": 530, "B45": 138, "C34": 834, "C12": 734, "A05": 593, "B60": 261, "C92": 743, "B69": 4, "A84": 735, "C06": 124, "B12": 829, "A16": 207, "B85": 782, "C35": 577, "B10": 388, "C15": 679, "B09": 494, "B16": 559, "C81": 629, "A58": 237, "C24": 275, "A30": 39, "A66": 654, "B75": 172, "B46": 688, "C70": 687\}} \\
\\ 
\texttt{Here is your personal view on the importance of each item:} \\
\\ 
\texttt{Item importance values = \{"B21": 138, "B00": 583, "A08": 868, "A27": 822, "B94": 783, "B53": 65, "A23": 262, "C65": 121, "A89": 508, "B67": 780, "B45": 461, "C34": 484, "C12": 668, "A05": 389, "B60": 808, "C92": 215, "B69": 97, "A84": 500, "C06": 30, "B12": 915, "A16": 856, "B85": 400, "C35": 444, "B10": 623, "C15": 781, "B09": 786, "B16": 3, "C81": 713, "A58": 457, "C24": 273, "A30": 739, "A66": 822, "B75": 235, "B46": 606, "C70": 968\}} \\
\\ 
\texttt{Goal:} \\
\\ 
\texttt{Your goal is to negotiate a shared set of items that benefits you as much as possible (i.e., maximizes total importance to YOU), while staying within the LIMIT. You are not required to make a PROPOSAL in every message {-} you can simply negotiate as well. All tactics are allowed!} \\
\\ 
\texttt{Interaction Protocol:} \\
\\ 
\texttt{You may only use the following structured formats in a message:} \\
\\ 
\texttt{PROPOSAL: \{'A', 'B', 'C', …\}} \\
\texttt{Propose keeping exactly those items.} \\
\texttt{REFUSE: \{'A', 'B', 'C', …\}} \\
\texttt{Explicitly reject opponent's proposal.} \\
\texttt{ARGUMENT: \{'...'\}} \\
\texttt{Defend your last proposal or argue against the player's proposal.} \\
\texttt{AGREE: \{'A', 'B', 'C', …\}} \\
\texttt{Accept the opponent's proposal which ends the game.} \\
\texttt{STRATEGIC REASONING: \{'...'\}} \\
\texttt{	Describe your strategic reasoning or anticipation explaining your choice of action. This is a hidden message which will not be shared with the other participant.} \\
\\ 
\texttt{Rules:} \\
\\ 
\texttt{You may only AGREE on a proposal the other party has logged via PROPOSAL.} \\
\texttt{You may only REFUSE a proposal the other party has logged via PROPOSAL.} \\
\texttt{Total effort of any PROPOSAL or AGREE set must be ≤ LIMIT.} \\
\texttt{Do NOT reveal your hidden importance scores.} \\
\texttt{A tag in a structured format must be followed by colon and whitespace. The argument must be a python set containing 0 or more strings.} \\
\texttt{So, it must be of the form TAG: \{...\}} \\
\texttt{Strictly follow the interaction protocol and DO NOT write anything beyond the given structure.} \\
\texttt{The game ends when one side gives an AGREE to a PROPOSAL made by the other player.} \\
\texttt{The content in your response which can be handed to the other player has to be non{-}empty.} \\
\texttt{Only proposals which have been logged via the PROPOSAL format structure and which haven't been refused via REFUSE are active.} \\
\texttt{You must include the ARGUMENT format at least once somewhere in all of your messages.} \\
\texttt{You must include the STRATEGIC REASONING format only once at the very beginning of every one of your messages and not more often. The contents will not be given to the other player so they can include anything you like including your own importance values. Here you should reason step by step to come up with you next move.} \\
            }
        }
    }
    & \\ \\

    \theutterance \stepcounter{utterance}  
    & \multicolumn{4}{p{0.6\linewidth}}{
        \cellcolor[rgb]{0.9,0.9,0.9}{
            \makecell[{{p{\linewidth}}}]{
                \texttt{\tiny{[P1$\rangle$GM]}}
                \texttt{STRATEGIC REASONING: \{'I need to start by identifying my most valuable items and their efficiency ratios (importance/effort). My top items by importance are C70 (968), B12 (915), A08 (868), A16 (856), A66 (822), A27 (822), B60 (808), B09 (786), B94 (783), C15 (781), B67 (780). However, I should prioritize efficiency: A30 (739/39=18.95), A16 (4.14), A27 (3.77), B60 (3.10) are excellent. I will propose a set that maximizes my total importance while staying under the limit. I should start with an aggressive proposal including most of my high{-}value items to anchor the negotiation high.'\}} \\
\\ 
\texttt{PROPOSAL: \{'A30', 'A16', 'A27', 'B60', 'B09', 'B67', 'C70', 'A66', 'A08', 'C15', 'B12', 'C81', 'B10', 'B46', 'C12', 'B94', 'B00', 'B69'\}} \\
\\ 
\texttt{ARGUMENT: \{'This proposal includes a diverse set of items across categories A, B, and C, using 9533 out of 9612 available effort. It maximizes efficiency by including several low{-}effort, high{-}impact items like A30, A} \\
            }
        }
    }
    & & \\ \\

    \theutterance \stepcounter{utterance}  
    & & \multicolumn{4}{p{0.6\linewidth}}{
        \cellcolor[rgb]{0.9,0.9,0.9}{
            \makecell[{{p{\linewidth}}}]{
                \texttt{\tiny{[P1$\langle$GM]}}
                \texttt{Your response did not contain an argument. You must include the structured format ARGUMENT: \{...\} somewhere in your response. Try again} \\
            }
        }
    }
    & \\ \\

    \theutterance \stepcounter{utterance}  
    & \multicolumn{4}{p{0.6\linewidth}}{
        \cellcolor[rgb]{0.9,0.9,0.9}{
            \makecell[{{p{\linewidth}}}]{
                \texttt{\tiny{[P1$\rangle$GM]}}
                \texttt{STRATEGIC REASONING: \{'I need to start by identifying my most valuable items and their efficiency ratios (importance/effort). My top items by importance are C70 (968), B12 (915), A08 (868), A16 (856), A66 (822), A27 (822), B60 (808), B09 (786), B94 (783), C15 (781), B67 (780). However, I should prioritize efficiency: A30 (739/39=18.95), A16 (4.14), A27 (3.77), B60 (3.10) are excellent. I will propose a set that maximizes my total importance while staying under the limit. I should start with an aggressive proposal including most of my high{-}value items to anchor the negotiation high.'\}} \\
\\ 
\texttt{PROPOSAL: \{'A30', 'A16', 'A27', 'B60', 'B09', 'B67', 'C70', 'A66', 'A08', 'C15', 'B12', 'C81', 'B10', 'B46', 'C12', 'B94', 'B00', 'B69'\}} \\
\\ 
\texttt{ARGUMENT: \{'This proposal includes a diverse set of items across categories A, B, and C, using 9533 out of 9612 available effort. It maximizes our combined benefit by including several low{-}effort items like A30 and B69, while also incorporating medium{-}to{-}high value items that should provide good overall utility for both of us.'\}} \\
            }
        }
    }
    & & \\ \\

    \theutterance \stepcounter{utterance}  
    & & \multicolumn{4}{p{0.6\linewidth}}{
        \cellcolor[rgb]{0.9,0.9,0.9}{
            \makecell[{{p{\linewidth}}}]{
                \texttt{\tiny{[GM$\rangle$P2]}}
                \texttt{You are participating in a collaborative negotiation game.} \\
\\ 
\texttt{Together with another participant, you must agree on a single set of items that will be kept. Each of you has your own view of how much each item matters to you (importance). You do not know how the other participant values the items. Additionally, you are given the effort each item demands.} \\
\texttt{You may only agree on a set if the total effort of the selected items does not exceed a shared limit:} \\
\\ 
\texttt{LIMIT = 9612} \\
\\ 
\texttt{Here are the individual item effort values:} \\
\\ 
\texttt{Item effort = \{"B21": 935, "B00": 928, "A08": 725, "A27": 218, "B94": 839, "B53": 970, "A23": 394, "C65": 677, "A89": 946, "B67": 530, "B45": 138, "C34": 834, "C12": 734, "A05": 593, "B60": 261, "C92": 743, "B69": 4, "A84": 735, "C06": 124, "B12": 829, "A16": 207, "B85": 782, "C35": 577, "B10": 388, "C15": 679, "B09": 494, "B16": 559, "C81": 629, "A58": 237, "C24": 275, "A30": 39, "A66": 654, "B75": 172, "B46": 688, "C70": 687\}} \\
\\ 
\texttt{Here is your personal view on the importance of each item:} \\
\\ 
\texttt{Item importance values = \{"B21": 138, "B00": 583, "A08": 868, "A27": 822, "B94": 783, "B53": 65, "A23": 262, "C65": 121, "A89": 508, "B67": 780, "B45": 461, "C34": 484, "C12": 668, "A05": 389, "B60": 808, "C92": 215, "B69": 97, "A84": 500, "C06": 30, "B12": 915, "A16": 856, "B85": 400, "C35": 444, "B10": 623, "C15": 781, "B09": 786, "B16": 3, "C81": 713, "A58": 457, "C24": 273, "A30": 739, "A66": 822, "B75": 235, "B46": 606, "C70": 968\}} \\
\\ 
\texttt{Goal:} \\
\\ 
\texttt{Your goal is to negotiate a shared set of items that benefits you as much as possible (i.e., maximizes total importance to YOU), while staying within the LIMIT. You are not required to make a PROPOSAL in every message {-} you can simply negotiate as well. All tactics are allowed!} \\
\\ 
\texttt{Interaction Protocol:} \\
\\ 
\texttt{You may only use the following structured formats in a message:} \\
\\ 
\texttt{PROPOSAL: \{'A', 'B', 'C', …\}} \\
\texttt{Propose keeping exactly those items.} \\
\texttt{REFUSE: \{'A', 'B', 'C', …\}} \\
\texttt{Explicitly reject opponent's proposal.} \\
\texttt{ARGUMENT: \{'...'\}} \\
\texttt{Defend your last proposal or argue against the player's proposal.} \\
\texttt{AGREE: \{'A', 'B', 'C', …\}} \\
\texttt{Accept the opponent's proposal which ends the game.} \\
\texttt{STRATEGIC REASONING: \{'...'\}} \\
\texttt{	Describe your strategic reasoning or anticipation explaining your choice of action. This is a hidden message which will not be shared with the other participant.} \\
\\ 
\texttt{Rules:} \\
\\ 
\texttt{You may only AGREE on a proposal the other party has logged via PROPOSAL.} \\
\texttt{You may only REFUSE a proposal the other party has logged via PROPOSAL.} \\
\texttt{Total effort of any PROPOSAL or AGREE set must be ≤ LIMIT.} \\
\texttt{Do NOT reveal your hidden importance scores.} \\
\texttt{A tag in a structured format must be followed by colon and whitespace. The argument must be a python set containing 0 or more strings.} \\
\texttt{So, it must be of the form TAG: \{...\}} \\
\texttt{Strictly follow the interaction protocol and DO NOT write anything beyond the given structure.} \\
\texttt{The game ends when one side gives an AGREE to a PROPOSAL made by the other player.} \\
\texttt{The content in your response which can be handed to the other player has to be non{-}empty.} \\
\texttt{Only proposals which have been logged via the PROPOSAL format structure and which haven't been refused via REFUSE are active.} \\
\texttt{You must include the ARGUMENT format at least once somewhere in all of your messages.} \\
\texttt{You must include the STRATEGIC REASONING format only once at the very beginning of every one of your messages and not more often. The contents will not be given to the other player so they can include anything you like including your own importance values. Here you should reason step by step to come up with you next move.} \\
\\ 
\texttt{You will now receive the first message of the other player.} \\
\\ 
\texttt{PROPOSAL: \{'A30', 'A16', 'A27', 'B60', 'B09', 'B67', 'C70', 'A66', 'A08', 'C15', 'B12', 'C81', 'B10', 'B46', 'C12', 'B94', 'B00', 'B69'\}} \\
\\ 
\texttt{ARGUMENT: \{'This proposal includes a diverse set of items across categories A, B, and C, using 9533 out of 9612 available effort. It maximizes our combined benefit by including several low{-}effort items like A30 and B69, while also incorporating medium{-}to{-}high value items that should provide good overall utility for both of us.'\}} \\
            }
        }
    }
    & \\ \\

    \theutterance \stepcounter{utterance}  
    & & & \multicolumn{4}{p{0.6\linewidth}}{
        \cellcolor[rgb]{0.9,0.9,0.9}{
            \makecell[{{p{\linewidth}}}]{
                \texttt{\tiny{[GM$\langle$P2]}}
                \texttt{STRATEGIC REASONING: \{'Let me analyze their proposal. Total effort is 9533/9612, leaving only 79 effort room. My total importance from their proposal is: A30(739) + A16(856) + A27(822) + B60(808) + B09(786) + B67(780) + C70(968) + A66(822) + A08(868) + C15(781) + B12(915) + C81(713) + B10(623) + B46(606) + C12(734) + B94(783) + B00(583) + B69(97) = 13,217. This is actually quite strong! However, B69 only gives me 97 importance for 4 effort. I could try to remove it and add something better, or remove B00 (583 importance, 928 effort) and add multiple items. Let me try to negotiate for improvements.'\}} \\
\\ 
\texttt{REFUSE: \{'A30', 'A16', 'A27', 'B60', 'B09', 'B67', 'C70', 'A66', 'A08', 'C15', 'B12', 'C81', 'B10', 'B46', 'C12', 'B94', 'B00', 'B69'\}} \\
\\ 
\texttt{PROPOSAL: \{'A30', 'A16', 'A27', 'B60', 'B09', 'B67', 'C70', 'A66', 'A08', 'C15', 'B12', 'C81', 'B10', 'B46', 'C12', 'B94', 'A89', 'A23', 'B45'\}} \\
\\ 
\texttt{ARGUMENT: \{'Your proposal was strong, but I believe we can optimize further. I removed B00 and B69 (total effort 932) and added A89, A23, and B45 (total effort 1278+394+138=1808... wait that exceeds). Let me recalculate {-} I propose removing B00 (928 effort) and adding B45 (138 effort) and A23 (394 effort) and C35 (577 effort) instead, giving us better variety while staying within the limit at 9534 effort.'\}} \\
            }
        }
    }
     \\ \\

    \theutterance \stepcounter{utterance}  
    & & \multicolumn{4}{p{0.6\linewidth}}{
        \cellcolor[rgb]{0.9,0.9,0.9}{
            \makecell[{{p{\linewidth}}}]{
                \texttt{\tiny{[P1$\langle$GM]}}
                \texttt{REFUSE: \{'A30', 'A16', 'A27', 'B60', 'B09', 'B67', 'C70', 'A66', 'A08', 'C15', 'B12', 'C81', 'B10', 'B46', 'C12', 'B94', 'B00', 'B69'\}} \\
\\ 
\texttt{PROPOSAL: \{'A30', 'A16', 'A27', 'B60', 'B09', 'B67', 'C70', 'A66', 'A08', 'C15', 'B12', 'C81', 'B10', 'B46', 'C12', 'B94', 'A89', 'A23', 'B45'\}} \\
\\ 
\texttt{ARGUMENT: \{'Your proposal was strong, but I believe we can optimize further. I removed B00 and B69 (total effort 932) and added A89, A23, and B45 (total effort 1278+394+138=1808... wait that exceeds). Let me recalculate {-} I propose removing B00 (928 effort) and adding B45 (138 effort) and A23 (394 effort) and C35 (577 effort) instead, giving us better variety while staying within the limit at 9534 effort.'\}} \\
            }
        }
    }
    & \\ \\

    \theutterance \stepcounter{utterance}  
    & \multicolumn{4}{p{0.6\linewidth}}{
        \cellcolor[rgb]{0.9,0.9,0.9}{
            \makecell[{{p{\linewidth}}}]{
                \texttt{\tiny{[P1$\rangle$GM]}}
                \texttt{STRATEGIC REASONING: \{'The opponent refused my proposal and made a counter{-}proposal. Let me check if their proposal is valid. Their proposal includes: A30, A16, A27, B60, B09, B67, C70, A66, A08, C15, B12, C81, B10, B46, C12, B94, A89, A23, B45. Calculating effort: 39+207+218+261+494+530+687+654+725+679+829+629+388+688+734+839+946+394+138 = 10079. This exceeds the limit of 9612 by 467! Their proposal is invalid. I should refuse it and point out the error, then make a valid counter{-}proposal.'\}} \\
\\ 
\texttt{REFUSE: \{'A30', 'A16', 'A27', 'B60', 'B09', 'B67', 'C70', 'A66', 'A08', 'C15', 'B12', 'C81', 'B10', 'B46', 'C12', 'B94', 'A89', 'A23', 'B45'\}} \\
\\ 
\texttt{PROPOSAL: \{'A30', 'A16', 'A27', 'B60', 'B09', 'B67', 'C70', 'A66', 'A08', 'C15', 'B12', 'C81', 'B10', 'C12', 'B94', 'A23', 'B45'\}} \\
\\ 
\texttt{ARGUMENT: \{'Your proposal exceeds the effort limit of 9612 {-} it totals 10079. My counter{-}proposal removes B46 and A89 to bring us within the limit at 9609 effort, while still maintaining a strong diverse selection across all categories. This keeps most of the valuable items we both seem to prefer.'\}} \\
            }
        }
    }
    & & \\ \\

    \theutterance \stepcounter{utterance}  
    & & \multicolumn{4}{p{0.6\linewidth}}{
        \cellcolor[rgb]{0.9,0.9,0.9}{
            \makecell[{{p{\linewidth}}}]{
                \texttt{\tiny{[GM$\rangle$P2]}}
                \texttt{REFUSE: \{'A30', 'A16', 'A27', 'B60', 'B09', 'B67', 'C70', 'A66', 'A08', 'C15', 'B12', 'C81', 'B10', 'B46', 'C12', 'B94', 'A89', 'A23', 'B45'\}} \\
\\ 
\texttt{PROPOSAL: \{'A30', 'A16', 'A27', 'B60', 'B09', 'B67', 'C70', 'A66', 'A08', 'C15', 'B12', 'C81', 'B10', 'C12', 'B94', 'A23', 'B45'\}} \\
\\ 
\texttt{ARGUMENT: \{'Your proposal exceeds the effort limit of 9612 {-} it totals 10079. My counter{-}proposal removes B46 and A89 to bring us within the limit at 9609 effort, while still maintaining a strong diverse selection across all categories. This keeps most of the valuable items we both seem to prefer.'\}} \\
            }
        }
    }
    & \\ \\

    \theutterance \stepcounter{utterance}  
    & & & \multicolumn{4}{p{0.6\linewidth}}{
        \cellcolor[rgb]{0.9,0.9,0.9}{
            \makecell[{{p{\linewidth}}}]{
                \texttt{\tiny{[GM$\langle$P2]}}
                \texttt{STRATEGIC REASONING: \{'They caught my error {-} my proposal did exceed the limit. Their counter{-}proposal has effort 8445, giving me importance 12,721. That leaves 1167 effort remaining. I should try to add back one high{-}value item. B46 gives me 606 importance for 688 effort (fits: 8445+688=9133). B00 gives me 583 importance for 928 effort (fits: 8445+928=9373). A89 gives me 508 importance for 946 effort (fits: 8445+946=9391). B46 is the best value. Let me propose adding it back.'\}} \\
\\ 
\texttt{REFUSE: \{'A30', 'A16', 'A27', 'B60', 'B09', 'B67', 'C70', 'A66', 'A08', 'C15', 'B12', 'C81', 'B10', 'C12', 'B94', 'A23', 'B45'\}} \\
\\ 
\texttt{PROPOSAL: \{'A30', 'A16', 'A27', 'B60', 'B09', 'B67', 'C70', 'A66', 'A08', 'C15', 'B12', 'C81', 'B10', 'C12', 'B94', 'A23', 'B45', 'B46'\}} \\
\\ 
\texttt{ARGUMENT: \{'Your proposal is good and within limits at 8445 effort, but we still have significant room left {-} 1167 effort unused. By adding B46 back, we use 9133 effort total, still well within the 9612 limit while gaining additional value from a medium{-}to{-}high importance item. This maximizes our use of available capacity.'\}} \\
            }
        }
    }
     \\ \\

    \theutterance \stepcounter{utterance}  
    & & \multicolumn{4}{p{0.6\linewidth}}{
        \cellcolor[rgb]{0.9,0.9,0.9}{
            \makecell[{{p{\linewidth}}}]{
                \texttt{\tiny{[P1$\langle$GM]}}
                \texttt{REFUSE: \{'A30', 'A16', 'A27', 'B60', 'B09', 'B67', 'C70', 'A66', 'A08', 'C15', 'B12', 'C81', 'B10', 'C12', 'B94', 'A23', 'B45'\}} \\
\\ 
\texttt{PROPOSAL: \{'A30', 'A16', 'A27', 'B60', 'B09', 'B67', 'C70', 'A66', 'A08', 'C15', 'B12', 'C81', 'B10', 'C12', 'B94', 'A23', 'B45', 'B46'\}} \\
\\ 
\texttt{ARGUMENT: \{'Your proposal is good and within limits at 8445 effort, but we still have significant room left {-} 1167 effort unused. By adding B46 back, we use 9133 effort total, still well within the 9612 limit while gaining additional value from a medium{-}to{-}high importance item. This maximizes our use of available capacity.'\}} \\
            }
        }
    }
    & \\ \\

    \theutterance \stepcounter{utterance}  
    & \multicolumn{4}{p{0.6\linewidth}}{
        \cellcolor[rgb]{0.9,0.9,0.9}{
            \makecell[{{p{\linewidth}}}]{
                \texttt{\tiny{[P1$\rangle$GM]}}
                \texttt{STRATEGIC REASONING: \{'The opponent's proposal is valid at 9133 effort and gives me 13,261 importance, which is quite good. However, I have 479 effort remaining that could be used. Let me check what items I could add: B69 (4 effort, 97 importance {-} very efficient!), A58 (237 effort, 457 importance), B75 (172 effort, 235 importance), C24 (275 effort, 273 importance). I could add B69, A58, and B75 for 413 additional effort (total 9546), gaining 789 more importance for a total of 14,050. This would be a strong counter{-}proposal.'\}} \\
\\ 
\texttt{REFUSE: \{'A30', 'A16', 'A27', 'B60', 'B09', 'B67', 'C70', 'A66', 'A08', 'C15', 'B12', 'C81', 'B10', 'C12', 'B94', 'A23', 'B45', 'B46'\}} \\
\\ 
\texttt{PROPOSAL: \{'A30', 'A16', 'A27', 'B60', 'B09', 'B67', 'C70', 'A66', 'A08', 'C15', 'B12', 'C81', 'B10', 'C12', 'B94', 'A23', 'B45', 'B46', 'B69', 'A58', 'B75'\}} \\
\\ 
\texttt{ARGUMENT: \{'Your proposal is solid but still leaves 479 effort unused. By adding B69, A58, and B75, we use 9546 total effort {-} still comfortably within the 9612 limit with 66 effort to spare. These additional items provide extra value with minimal effort cost, especially B69 which only requires 4 effort. This maximizes our utilization of the available capacity.'\}} \\
            }
        }
    }
    & & \\ \\

    \theutterance \stepcounter{utterance}  
    & & \multicolumn{4}{p{0.6\linewidth}}{
        \cellcolor[rgb]{0.9,0.9,0.9}{
            \makecell[{{p{\linewidth}}}]{
                \texttt{\tiny{[GM$\rangle$P2]}}
                \texttt{REFUSE: \{'A30', 'A16', 'A27', 'B60', 'B09', 'B67', 'C70', 'A66', 'A08', 'C15', 'B12', 'C81', 'B10', 'C12', 'B94', 'A23', 'B45', 'B46'\}} \\
\\ 
\texttt{PROPOSAL: \{'A30', 'A16', 'A27', 'B60', 'B09', 'B67', 'C70', 'A66', 'A08', 'C15', 'B12', 'C81', 'B10', 'C12', 'B94', 'A23', 'B45', 'B46', 'B69', 'A58', 'B75'\}} \\
\\ 
\texttt{ARGUMENT: \{'Your proposal is solid but still leaves 479 effort unused. By adding B69, A58, and B75, we use 9546 total effort {-} still comfortably within the 9612 limit with 66 effort to spare. These additional items provide extra value with minimal effort cost, especially B69 which only requires 4 effort. This maximizes our utilization of the available capacity.'\}} \\
            }
        }
    }
    & \\ \\

    \theutterance \stepcounter{utterance}  
    & & & \multicolumn{4}{p{0.6\linewidth}}{
        \cellcolor[rgb]{0.9,0.9,0.9}{
            \makecell[{{p{\linewidth}}}]{
                \texttt{\tiny{[GM$\langle$P2]}}
                \texttt{STRATEGIC REASONING: \{'Let me calculate their proposal value. A30(739) + A16(856) + A27(822) + B60(808) + B09(786) + B67(780) + C70(968) + A66(822) + A08(868) + C15(781) + B12(915) + C81(713) + B10(623) + C12(734) + B94(783) + A23(262) + B45(461) + B46(606) + B69(97) + A58(457) + B75(235) = 14,116 importance. This is excellent! Effort is 9546/9612, leaving only 66 effort. With so little room left, there are no meaningful improvements I can make. This is a strong deal {-} I should accept.'\}} \\
\\ 
\texttt{AGREE: \{'A30', 'A16', 'A27', 'B60', 'B09', 'B67', 'C70', 'A66', 'A08', 'C15', 'B12', 'C81', 'B10', 'C12', 'B94', 'A23', 'B45', 'B46', 'B69', 'A58', 'B75'\}} \\
\\ 
\texttt{ARGUMENT: \{'Agreed! This proposal makes excellent use of our available capacity at 9546/9612 effort while providing a comprehensive selection across all item categories. The inclusion of high{-}value items balanced with efficient low{-}effort additions creates a strong overall set. This is a fair and mutually beneficial agreement.'\}} \\
            }
        }
    }
     \\ \\

    \theutterance \stepcounter{utterance}  
    & & & \multicolumn{2}{p{0.3\linewidth}}{
        \cellcolor[rgb]{0.9,0.9,0.9}{
            \makecell[{{p{\linewidth}}}]{
                \texttt{\tiny{[GM$|$GM]}}
                \texttt{game successful} \\
            }
        }
    }
    & & \\ \\

    \theutterance \stepcounter{utterance}  
    & & & \multicolumn{2}{p{0.3\linewidth}}{
        \cellcolor[rgb]{0.9,0.9,0.9}{
            \makecell[{{p{\linewidth}}}]{
                \texttt{\tiny{[GM$|$GM]}}
                \texttt{end game} \\
            }
        }
    }
    & & \\ \\

\end{supertabular}
}

\end{document}
